\section{Introduction}

\subsection{The Problem and Its Urgency}

Credit card fraud has become one of the fastest-growing financial crimes worldwide.
The Nilson Report estimates global card fraud losses at \textbf{USD 27.33 billion in
2023}, with projections exceeding USD 40 billion by 2027 \cite{Nilson2024}. The
problem is not only the direct loss. A 2024 LexisNexis study establishes that
every \textbf{USD 1 lost to fraud costs financial institutions USD 4.45} to detect,
prevent and compensate \cite{LexisNexis2024}. Beyond the direct financial damage,
fraud has a corrosive effect on customer trust: a substantial share of fraud
victims disengage from their bank within months of the incident.

Banks respond with two operational levers that both carry hidden costs. They
\emph{reject} suspicious transactions, which avoids fraud losses but punishes
legitimate customers. They \emph{manually review} borderline cases, which is
accurate but consumes scarce human analyst time. Traditional rule-based systems
--- for example, "block every foreign transaction above \$X" --- are cheap to
implement but are blunt instruments: they miss a large share of fraud while
wrongly blocking ordinary customers. The core engineering problem is therefore
not simply ``detect fraud'' but \emph{decide, per transaction, how confident we
are, and what action that confidence justifies}.

\subsection{Research Motivation}

This project applies a supervised learning pipeline to the canonical credit card
fraud dataset \cite{ULB2013}, but the evaluation philosophy is deliberately
business-facing rather than metric-facing. Three commitments follow from this:

\begin{enumerate}
    \item \textbf{Accuracy is not a meaningful metric on this problem.} With a
    class ratio of roughly 1:599, a classifier that labels \emph{every}
    transaction legitimate scores 99.83\% accuracy while catching exactly zero
    frauds. We therefore evaluate Recall, Precision, F1 and Average Precision.
    \item \textbf{Imbalance handling must be justified economically.} SMOTE
    \cite{Chawla2002}, Random Undersampling, SMOTE-ENN
    \cite{Guo2008} and class weighting are all routinely prescribed. We compare
    all of them and translate each into a monetary cost, because a technique that
    improves a metric but explodes the analyst workload is worthless to a bank.
    \item \textbf{The output must be an operating decision, not a probability.} A
    bank needs to know which transactions to block automatically, which to send
    to an analyst, and which to approve --- and what that policy costs per
    million transactions.
\end{enumerate}

\subsection{Research Questions}

The project is organised around four questions, each of which is answered
by a specific part of this report:

\begin{enumerate}
    \item \textbf{RQ1 (Detection).} Which model family and imbalance-handling
    technique maximise fraud detection on this dataset under a leakage-controlled
    evaluation protocol? (Sections 3--4.)
    \item \textbf{RQ2 (Cost).} At a precision floor of 95\%, what does each
    operating point cost in monetary terms, and how does the ranking of
    techniques change once a business cost is attached to a false alarm?
    (Section 5.)
    \item \textbf{RQ3 (Capacity).} What is the maximum fraud that can be
    caught if the bank can review only 100 flagged transactions per day, and
    what decision architecture satisfies that constraint? (Section 5.)
    \item \textbf{RQ4 (Validity).} How much of the measured performance
    survives the threats to validity that this design admits --- threshold
    leakage, model selection on the test set, a two-day sample and a single
    random seed? (Sections 4.5 and 6.2.)
\end{enumerate}

\subsection{Problem Definition}

\begin{itemize}
    \item \textbf{Input:} One credit card transaction, represented by 30 original
    numerical attributes (\texttt{Time} in seconds since the first transaction,
    \texttt{V1}--\texttt{V28} obtained by PCA, and \texttt{Amount} in USD),
    extended by five engineered features (\texttt{Hour\_of\_Day}, \texttt{Day},
    \texttt{Is\_Night}, \texttt{Amount\_log}, \texttt{Is\_Zero\_Amount}).
    \item \textbf{Output:} A fraud probability, converted by a decision
    threshold into one of three operational actions: \emph{auto-block},
    \emph{manual review}, or \emph{approve}.
    \item \textbf{Application:} Real-time transaction authorisation at the
    issuing bank, as one decision stage inside the wider card-network risk
    process. We model the decision itself --- what to score, at what cost, and
    who acts on the output --- and make no claim about any specific vendor
    product.
\end{itemize}

\subsection{Literature Review}

\paragraph*{Imbalanced data as the central issue.}
Dal Pozzolo et al. \cite{DalPozzolo2014} established that ``imbalanced data is
the central issue in credit card fraud detection'' and combined SMOTE with a
pairwise ranking learner, introducing calibration as a means of recovering
probabilities under heavy resampling. Their contribution remains the standard
baseline: SMOTE, evaluated with Average Precision rather than accuracy, and
probabilities that can be thresholded meaningfully.

\paragraph*{Comparative benchmarking.}
West and Bhattacharya \cite{West2016} surveyed statistical methods for fraud
detection and emphasised the gap between laboratory accuracy and operational
performance. Leal-Avila et al. \cite{LealAvila2023} benchmarked Random Forest
\cite{Breiman2001}, XGBoost \cite{Jian2013} and Isolation Forest on exactly the
ULB dataset used here, confirming that tree ensembles are state of the art on
this data while unsupervised anomaly detection is materially weaker.

\paragraph*{Streaming and real-time deployment.}
Carcillo et al. \cite{Carcillo2018} framed the problem as a streaming task and
built SCARFF on Apache Spark, making the operational latency requirement
explicit: scoring must complete within the authorisation window, not overnight.

\paragraph*{The role of calibration.}
Van Calster et al. \cite{VanCalster2019} showed that a model may rank
transactions perfectly yet be badly calibrated, which matters directly when a
threshold is chosen in business terms. Davis and Goadrich \cite{Davis2006}
established the theoretical relationship between the precision--recall and ROC
curves, justifying our choice of Average Precision as the primary metric on an
imbalanced dataset \cite{Fawcett2006}.

\paragraph*{Explainability as a requirement, not a luxury.}
Under regulatory pressure for explainable AI, Lundberg and Lee \cite{Lundberg2017}
introduced SHAP, with TreeSHAP \cite{Lundberg2020} providing exact Shapley
values for tree ensembles. A bank must be able to answer ``why was this
transaction blocked?'' to the customer and to the regulator.

\paragraph*{The gap addressed by this project.}
Existing work typically reports either predictive metrics or operational
observations, rarely both in a single reproducible pipeline, and almost never
converts the model output into a costed, three-tier operating policy evaluated
against explicit rule-based baselines. This project connects all three: a
reproducible pipeline, a rigorous model comparison across imbalance techniques,
and a costed decision policy benchmarked against the rules a bank would
otherwise deploy.

\subsection{Contribution}

Our contributions are: (i) an empirical, leakage-controlled comparison of five
class-imbalance techniques across five model families on the ULB dataset; (ii) an
observation that a wider hyper-parameter search \emph{degrades} test
performance when only 331 fraud examples are available --- an observation we
report with its statistical caveat rather than as a proof; (iii) a quantified
assessment of the cost of each imbalance technique, including the mean-versus-best
distinction that the technique ranking depends on; and (iv) a two-tier operating
policy that attains higher recall than any single threshold that respects the
100-per-day review capacity stated in the brief, while requiring only 70 analyst
reviews per day.

\subsection{Contribution of Group Members}
Table~\ref{tab:contrib} records the individual contribution of each group
member, documented so that it can be verified against the artefacts described
in Appendix~\ref{app:repro}.

\begin{table}[H]
\centering
\caption{Individual contributions of the group members.}
\label{tab:contrib}
\begin{tabular}{p{4.0cm}p{10.6cm}}
\toprule
\textbf{Member} & \textbf{Contribution} \\
\midrule
Nguyễn Tiến Công &
Introduction, problem definition, research motivation, literature review and
related work (Sections 1--2). \\
\addlinespace
Nguyễn Tuấn Sỹ &
Methodology and reproducibility (Section 3 and Appendix~A): pipeline
architecture, preprocessing, feature engineering, imbalance techniques,
hyper-parameter search, leakage control. \\
\addlinespace
Nguyễn Nho Thắng &
Results (Section 4): data quality, the 27-configuration imbalance benchmark,
model selection, threshold-selection leakage, explainability. \\
\addlinespace
Nguyễn Quang Cường &
Cost analysis and business decision (Section 5): cost model, analyst capacity,
the two-tier operating policy, scalability, the seven business questions. \\
\addlinespace
Nguyễn Gia Phong &
Discussion, limitations, recommendations and conclusion (Sections 6--7), plus
delivery of the presentation. \\
\bottomrule
\end{tabular}
\end{table}

\section{Related Work and Background}
\label{sec:background}

\subsection{Class Imbalance and Resampling}

The imbalance ratio here is 1:599, well beyond the point at which standard
cross-entropy training effectively ignores the minority class \cite{Guo2008}.
SMOTE \cite{Chawla2002} synthesises minority points by interpolating between
k-nearest minority neighbours, preserving the local geometry of the fraud
cluster. Random Undersampling instead discards majority examples, which is cheap
but destroys information and, as we demonstrate empirically, is the most
damaging of the techniques we test. SMOTE-ENN combines both and then removes
samples whose nearest neighbours are of mixed class.

\subsection{Model Families}

Logistic Regression serves as the linear baseline. Decision Tree provides
complete interpretability but high variance. Random Forest \cite{Breiman2001}
averages deep trees on bootstrap samples, making it robust to the heavy-tailed
amount distribution. XGBoost \cite{Jian2013} fits trees greedily with second-order
information and regularisation, typically the strongest tabular learner.
Isolation Forest detects anomalies by measuring the average path length of
isolated points and requires no labels, which we use to test whether fraud is
detectable as a pure anomaly.

\subsection{Metrics and Evaluation}

For a $1:599$ ratio, Accuracy is uninformative and AUC-ROC is dominated by the
negative class \cite{Fawcett2006}. Average Precision --- the area under the
precision--recall curve --- is the appropriate summary statistic
\cite{Davis2006}. Because the operating threshold is a business decision, we
report the full precision--recall curve and select thresholds subject to a
precision constraint of 95\%, corresponding to blocking at most five in every
one hundred flagged transactions.

\subsection{Cost Model}

We adopt a transparent cost model. A missed fraud costs
\textbf{USD 5.45} (USD 1.00 direct loss plus USD 4.45 handling, per
\cite{LexisNexis2024}). A false alarm costs \textbf{USD 0.30}, our own
assumption for a manual review. Total cost is therefore
\[
    C = n_{FN} \cdot 5.45 \;+\; n_{FP} \cdot 0.30 .
\]
Every conclusion in Section~\ref{sec:cost} is expressed under this model, and we
explicitly state in Section~\ref{sec:discussion} which conclusions depend on it.
